\begin{helper}
Let $k$ be a natural number and let $w_{ab}$ be nonnegative real weights.
Use the symmetric unordered-pair convention
\[
e(a,b)=
\begin{cases}
w_{ab},&a<b,\\
w_{ba},&b<a,\\
0,&a=b.
\end{cases}
\]
Let
\[
M=\sum_{a<b} e(a,b).
\]
Then, for every vertex $i$,
\[
0\le \sum_b e(i,b)
\quad\text{and}\quad
\sum_b e(i,b)\le M,
\]
and the incident totals satisfy
\[
\sum_i\sum_b e(i,b)=2M.
\]
\end{helper}
\begin{proof}
The symmetric convention and nonnegativity of the original weights give
$e(a,b)\ge 0$ for all $a,b$, $e(i,i)=0$, and $e(a,b)=e(b,a)$.
Therefore each incident total $\sum_b e(i,b)$ is a sum of nonnegative
terms, so it is nonnegative.

Fix $i$.  The terms of $\sum_b e(i,b)$ are exactly the weights of the
increasing unordered pairs incident with $i$: if $i<b$, the contributing
pair is $(i,b)$ and the weight is $e(i,b)$; if $b<i$, the contributing
pair is $(b,i)$ and its weight is $e(b,i)=e(i,b)$; and the diagonal term
$b=i$ is zero.  Hence $\sum_b e(i,b)$ is a sub-sum of
$M=\sum_{a<b}e(a,b)$, and all omitted terms are nonnegative.  Thus
$\sum_b e(i,b)\le M$.

Finally apply \noderef{KPartiteIncidentWeightDoubleCount} to the original
weights $w_{ab}$.  Its displayed symmetric incident sum is precisely
$\sum_i\sum_b e(i,b)$ under the definition of $e$, and its increasing-pair
sum is precisely $M$.  Therefore
$\sum_i\sum_b e(i,b)=2M$.
\end{proof}
